\chapter{Validation by propagation}
\label{ch:validation}

A fitted model is judged by whether it \emph{predicts}: whether it can be run forward from the recorded thrust
alone and stay close to what the vehicle actually did. \file{validation/propagate_model.py}
does this offline over any recorded bag and writes the result into a new MCAP file for
Foxglove.

\section{Algorithm}

\begin{lstlisting}[caption={Outline of \code{propagate_model.main()} and \code{propagate()}.},label={lst:propagate}]
model = load_model(path)              # refuse if m, X_udot, X_u, X_uu are null
odom, servo = read_tracks(bag)        # EKF on header.stamp, servo on log_time
servo.log_t -= thrust.servo_delay_s
verify_timestamps(odom, servo)        # rates, gaps, clock offset, overlap
sign_check(pwm, ekf_surge)            # does PWM > 1500 mean forward motion?
for each EKF sample k:                # output grid ~30 Hz
    for each sub-step of <= dt:       # internal grid, default 10 ms
        T   = thrust(held_pwm, held_voltage)
        psi, v = interp(EKF yaw, EKF sway)
        x, y  += h * R(psi) @ [u, v]   # eq. (planar kinematics)
        u      = rk4_step(u, T, h)     # eq. (surge model)
    if pwm stale:        restart from EKF (thrust unknown, not zero)
    if horizon >= reinit_s: restart from EKF (x, y, u)
    z = EKF z                          # heave is measured
write augmented bag + summary
\end{lstlisting}

The sign check correlates $\mathrm{sign}(p - p_0)$ with the sign of the EKF surge velocity. A
negative correlation is \emph{reported} with the suggestion \code{thrust.sign: -1}, but never
applied automatically.

\section{Free run versus re-initialisation}

\paragraph{Free run} (\code{--reinit-s 0}, or the default \qty{1800}{\second}, which is longer than any bag).
The model is started once from the EKF and never corrected. A wrong thrust scale, drag or
mass shows up as steadily growing error that nothing can hide. This is the verdict on whether
the model is right.

\paragraph{Re-initialisation} (e.g.\ \code{--reinit-s 10}). Every $N$ seconds position and surge speed
are reset to the EKF. This measures a narrower property, how well the model predicts $N$
seconds ahead from a known state, and it yields many independent samples instead of one. The
\emph{shape} of the error against the time since the reset (\code{/sysid/horizon_s}) is a
diagnostic of which parameter is wrong.

\section{Reading the error-versus-horizon curve}

\paragraph{Thrust scale or drag wrong.} Both shift the steady-state speed \eqref{eq:uss}. The
model then has a persistent velocity bias $\delta u$ and the position error grows
\emph{linearly}, $\delta x \approx \delta u\,t$.

\paragraph{Inertia wrong.} $\mtot$ does not affect the steady state. It only changes how fast
each transient settles. After a thrust change the velocity error is a transient that decays
with $\taueff$, so the position error from each transient is \emph{bounded} (of order
$\delta\taueff\,\Delta u$). For horizons much shorter than $\taueff$ the error starts out
quadratic ($\tfrac12\delta a\,t^2$, the case described in \file{maneuvers.md}). Over longer
horizons it grows only as successive transients add up, with random signs, so roughly like
$\sqrt t$.

\begin{figure}[htbp]
\centering
\begin{tikzpicture}
\begin{axis}[wide,xlabel={Horizon since re-initialisation / \unit{\second}},ylabel={RMS position error / \unit{\meter}},xmin=0,xmax=10,ymin=0]
\addplot[red] table[x=h,y=thrust]{data/horizon.dat};
\addplot[blue,dashed] table[x=h,y=mass]{data/horizon.dat};
\legend{thrust \qty{15}{\percent} low,$\mtot$ \qty{50}{\percent} high}
\end{axis}
\end{tikzpicture}
\caption{RMS position error against prediction horizon for a surge model whose thrust is
\qty{15}{\percent} too low and for one whose total inertia is \qty{50}{\percent} too high,
averaged over \qty{10}{\second} prediction windows started every \qty{2}{\second} on
\qty{600}{\second} of a simulated
piecewise-constant thrust profile, with the dummy BlueROV2 coefficients as ground truth.}
\label{fig:horizon}
\end{figure}
\FloatBarrier

\Cref{fig:horizon} shows the two signatures. The thrust error grows almost linearly. The
inertia error, even at \qty{50}{\percent}, grows sub-linearly and stays small. With the dummy
coefficients $\taueff$ is a few tenths of a second, so the quadratic phase is too short to see,
and the practical signature of an inertia error is a \emph{saturating}, step-correlated error
rather than a quadratic one. In practice: a straight line indicates a check of thrust and drag, and an
error concentrated just after thrust changes indicates a check of inertia and delay.

\section{Metrics in the summary}

For each predictor (model, and the power-law baseline) the summary reports:
\begin{itemize}
  \item \textbf{Surge RMS error and its ratio to the EKF surge standard deviation.} A ratio of
        \qty{100}{\percent} means the prediction is no better than always guessing the mean speed.
        For an unbiased predictor the ratio is $\sqrt{1 - R^2}$.
  \item \textbf{Position RMS error and final drift}, absolute and as a percentage of the distance travelled.
  \item \textbf{Path length}, predicted against EKF. This is a quick check of the overall speed scale.
\end{itemize}
The baseline matters because a dynamic model that cannot beat \eqref{eq:powerlaw} has not
earned its complexity.

\section{Output topics}

\begin{table}[htbp]
\centering
\caption{Topics added to the augmented bag by \file{propagate_model.py}.}
\label{tab:topics}
\small
\begin{tabularx}{\textwidth}{@{}l l X@{}}
\toprule
Topic & Type & Content \\
\midrule
\code{/sysid/model_odom} & \code{Odometry} & Model prediction, same frames as the EKF \\
\code{/sysid/powerlaw_odom} & \code{Odometry} & Baseline prediction \\
\code{/sysid/model_navsatfix} & \code{NavSatFix} & Prediction projected to latitude and longitude for the Map panel \\
\code{/sysid/model_info} & \code{String} & Version tag, coefficients, settings (\qty{1}{\hertz}) \\
\code{/sysid/error/position_m} & \code{Float64} & Horizontal position error against the EKF \\
\code{/sysid/error/surge_mps} & \code{Float64} & Surge speed error against the EKF \\
\code{/sysid/horizon_s} & \code{Float64} & Time since the last re-initialisation \\
\code{/sysid/thruster_forces} & \code{Float64MultiArray} & Force on all six channels (NaN = no curve) \\
\code{/sysid/battery_voltage} & \code{Float64} & Voltage the thrust was evaluated at \\
\bottomrule
\end{tabularx}
\end{table}
\FloatBarrier

\section{Geographic projection: the Kabsch fit}

The \code{odom} frame is local. To draw the prediction on a map, the transform from
\code{odom} to UTM is needed. \code{/gps/filtered} is the EKF output pushed through
\code{navsat_transform}'s fixed transform, so a rigid 2-D transform $\vect{y}_k = \vect{R}\vect{x}_k + \vect{t}$
fits it exactly. \code{fit_geo_anchor()} recovers it with the Kabsch algorithm \cite{kabsch1976}:
\begin{enumerate}
  \item centre both point sets: $\tilde{\vect{X}} = \vect{X} - \bar{\vect{x}}$, $\tilde{\vect{Y}} = \vect{Y} - \bar{\vect{y}}$;
  \item take the SVD $\tilde{\vect{X}}^\top\tilde{\vect{Y}} = \vect{U}\vect{\Sigma}\vect{V}^\top$;
  \item $\vect{R} = \vect{V}\,\mathrm{diag}(1, d)\,\vect{U}^\top$ with
        $d = \mathrm{sign}\det(\vect{V}\vect{U}^\top)$, which guarantees a rotation rather than a reflection;
  \item $\vect{t} = \bar{\vect{y}} - \vect{R}\bar{\vect{x}}$.
\end{enumerate}
The scale is pinned to 1 because metres in \code{odom} are metres in UTM. The fit is valuable
because the map frame is \emph{not} ENU-aligned: it is off by about \qty{17}{\degree} on the June bags.
Assuming ENU would put the predicted track in the wrong place.

\section{Reading the existing summaries}

\file{recordings/sysid/} holds several summaries produced while the pipeline was being built.
All of them used the dummy BlueROV2 model, so none says anything about POLARIS
coefficients. They are still instructive for how the outputs are read.

\begin{table}[htbp]
\centering
\caption{Key figures of the propagation summaries in \file{recordings/sysid/} (dummy
BlueROV2 model, free run).}
\label{tab:summaries}
\small
\begin{tabular}{@{}l l S[table-format=1.3] S[table-format=3.1] S[table-format=4.1] S[table-format=4.1]@{}}
\toprule
Bag & Predictor & {Surge RMS / \unit{\meter\per\second}} & {Ratio / \unit{\percent}} & {Path pred. / \unit{\meter}} & {Path EKF / \unit{\meter}} \\
\midrule
\code{zermatt_rectangle_07} & model & 0.089 & 76.6 & 2.8 & 9.1 \\
\code{zermatt_rectangle_07} & power law & 0.093 & 79.5 & 2.9 & 9.1 \\
\code{aut_..._grid_03} & model & 1.198 & 100.0 & 298.2 & 2718.8 \\
\code{aut_..._grid_03} & power law & 1.198 & 100.1 & 329.9 & 2718.8 \\
\bottomrule
\end{tabular}
\end{table}
\FloatBarrier

\paragraph{\code{zermatt_rectangle_07}.} Both predictors explain only about
\qty{40}{\percent} of the surge variance ($1 - 0.766^2$) and cover about a third of the EKF path. The model and the power law fail in the same way
even though they have nothing in common except the PWM input. That points at the input rather than at
either model. The servo stream in the April rectangle bags runs at only about \qty{2}{\hertz} and the
surge PWM never exceeds \qty{1683}{\micro\second}, so much of the motion may be driven by
something the surge channel does not see (other thrusters, current, sway).

\paragraph{\code{aut_no_z_surface_big_40_grid_03}.} A ratio of exactly \qty{100}{\percent} for
\emph{both} predictors, and an EKF path nine times longer than either prediction, are the
signature of a broken reference rather than a bad model. A direct check of the bag confirms it:
the EKF surge velocity reaches \qty{62}{\meter\per\second} in places (182 samples above
\qty{2}{\meter\per\second}), and the EKF position jumps by up to \qty{842}{\meter} between
consecutive samples. Without jumps larger than \qty{1}{\meter} the EKF path is \qty{479}{\meter},
not \qty{2719}{\meter}. The spikes inflate the ``spread'' to \qty{1.2}{\meter\per\second} and
dominate every metric. This bag should be used only on cropped windows where the EKF is
continuous.

\begin{keyidea}
Validation compares the model against the EKF, and the EKF is also only an estimate. Before a
large error is blamed on the model, the reference has to be checked: velocity spikes, position jumps,
resets. Agreement between model and baseline on a bad number is a strong hint that the shared
input or the reference is at fault.
\end{keyidea}

\begin{questions}
\begin{enumerate}
  \item Why is a free-run error the ``verdict'' but a re-initialised error only a ``diagnostic''?
  \item The position error grows linearly with horizon at \qty{0.05}{\meter\per\second}. Which
        coefficients are suspects, and which manoeuvre would separate them?
  \item What would \file{propagate_model.py} report if the surge thruster were mounted backwards?
\end{enumerate}
\end{questions}
